% ================================================================
\section{Introduzione}

Questo lavoro descrive il programma PythonGranularEngine (PGE) per la sintesi granulare in
tempo differito. La sintesi granulare costruisce il suono come
collezione di migliaia di grani, brevi eventi sonori di durata tipica
1-50~ms, e richiede di specificare, direttamente o per via
algoritmica, i parametri di ciascuno. Già la prima implementazione
computer documentata~\cite{Roads1978} osserva che la specifica
esplicita diventa intrattabile non appena la densità supera poche
unità per secondo. Il problema centrale del sistema è quindi il
controllo parametrico: rendere \emph{esplicito} e \emph{leggibile} il
governo di una massa di grani che nessun compositore può razionalizzare
grano per grano.

Una precisazione tassonomica è dovuta fin da subito. \emph{Sintesi
granulare} è qui termine-ombrello per il paradigma; il sistema, nello
specifico, non sintetizza grani da una primitiva di Gabor ma
\emph{granula} materiale registrato: legge e ricompone frammenti di
un campione sorgente. Appartiene cioè al ramo che
Lippe~\cite{Lippe1993cim} distingue come \emph{granular sampling}, la
granulazione di campioni (o micromontage), in cui il parametro
espressivo dominante è la posizione di lettura nel materiale. Questa
collocazione, che riprendo in \S\ref{sec:partitura} e
\S\ref{sec:tradizione}, motiva sia l'asse della partitura sia il
ruolo centrale del \texttt{PointerController}.

Il sistema adotta il tempo differito oggi in una fase storico-informatica in cui il tempo reale 
è tecnicamente disponibile; sulle ragioni di questa scelta
discuto in \S\ref{sec:implicazioni}. L'esposizione muove dal problema
concreto del controllo parametrico e sviluppa tre nuclei principali:
\begin{itemize}
\item optando per un approccio dichiarativo al controllo dei parametri granulari,
  un \emph{domain-specific language}%
  \footnote{Con \emph{domain-specific language} si intende un linguaggio 
  formale progettato per un dominio specifico, in contrapposizione ai 
  linguaggi general-purpose. Non richiede necessariamente una nuova sintassi e
  può appoggiarsi a un formato esistente come \emph{syntax host}, 
arricchendolo di un vocabolario e di una semantica propri del dominio. 
  }
  (DSL), validato da un \emph{Language Server}\footnote{Un \emph{Language Server} è un componente software che analizza 
il testo di un file durante la scrittura e fornisce in tempo reale 
autocompletamento, validazione e documentazione contestuale all'interno 
dell'editor. Il protocollo standard è LSP (\emph{Language Server Protocol}, 
Microsoft 2016).} (orientato alla 
  composizione assistita) durante
  la scrittura e interpretato come \emph{tendency mask} (traiettoria
  centrale, range di deviazione, campionamento per-grano). In particolare
  il nostro DSL sarà qui YAML\footnote{Acronimo ricorsivo di YAML Ain't 
  Markup Language, è un formato di serializzazione dati facilmente 
  leggibile dagli esseri umani.};
\item un sistema che differenzia le voci lungo
  quattro assi ortogonali (\emph{pitch}, \emph{onset}, \emph{pointer}, \emph{pan}) mediante
  strategie intercambiabili (accordo, lineare, stocastico, \dots) e un
  parametro \texttt{scatter} che controlla la variazione di asincronia temporale delle voci;
\item una partitura grafica generata automaticamente, il cui asse
  verticale codifica la posizione di lettura nel buffer sorgente,
  e un workflow di stream audio separati con cache incrementale ed export verso DAW (REAPER),
  che chiude il ciclo di generazione compositiva e verifica.
\end{itemize}
Il paper procede dal basso: prima il sistema (\S\ref{sec:architettura})
e la sua partitura (\S\ref{sec:partitura}), poi il posizionamento
nella tradizione granulare (\S\ref{sec:tradizione}), infine le
implicazioni teorico-compositive (\S\ref{sec:implicazioni}).




